\section{Baseline Recipe and Evaluation Protocol}
\label{sec:phase1}

The baseline stage adapts the base paper's KAN-ODE recipe
\cite{koenig2024kanode} (differences are listed in
Table~\ref{tab:paper-audit}) on the Lotka-Volterra predator-prey system
(Equation~\ref{eq:lv}), using a two-layer Kolmogorov-Arnold Network,
$[2,10,2]$, with $G=5$ Gaussian radial-basis-function grid points per edge
(240 parameters), trained on $t\in[0,3.5]$ (36 points, almost exactly one
predator-prey cycle) and scored strictly on $t\in(3.5,14.0]$ (105 unseen
points, about 3.2 further cycles). Extrapolation is scored this way, on
points the training loss never touched, because it is the only meaningful
test of whether the network learned the underlying dynamics rather than
memorized the shape of the training window. A model that only reproduces
the 36 training points could still be badly wrong about the system it is
meant to describe.

Two additional systems were carried forward alongside Lotka-Volterra: the
SIR epidemic model (Equation~\ref{eq:sir}) and the damped pendulum
(Equation~\ref{eq:pendulum}), with the splits of
Table~\ref{tab:datasets}. Both were chosen because their governing
equations are textbook-simple and well understood, so any training failure
can be attributed to the architecture or the optimization process rather
than to ambiguous target dynamics. Each also carries a physical law that is
separate from trajectory error. SIR conserves $S+I+R=1$ exactly, so any
model that violates it has learned a vector field whose structure is wrong,
not merely inaccurate. The pendulum's energy must decrease monotonically,
which a model can violate while still fitting the trajectory closely.

\subsection{Evaluation protocol}

This stage established the evaluation protocol used throughout the rest of
this report. Model selection uses the minimum \emph{training}-window loss
only (Algorithm~\ref{alg:train}). The headline error is the mean squared
error (MSE, Equation~\ref{eq:metric-mse}) between the model's predicted
trajectory and the true trajectory, computed strictly outside the training
horizon, so that a number cannot improve simply by widening what counts as
fit. The coefficient of determination, $R^2$ (Equation~\ref{eq:metric-r2}),
reports the same comparison on a scale where $1$ is a perfect match and $0$
is no better than predicting the mean of the true trajectory; a negative
value means the model performs worse than that trivial baseline. Every
later study keeps this selection rule, and says so where it reports a
different quantity (for example the full-horizon MSE on the pendulum).

\FloatBarrier
